\section{Command reference}
\label{sec:command-reference}

\begin{lstlisting}
# --- pool -----------------------------------------------------------
gpu status
gpu claim 1
gpu shell                       # interactive
gpu shell /path/to/script.sh    # one command
gpu release

# --- inside the container -------------------------------------------
nvidia-smi
ps -eo pid,etime,cmd | grep "[b]in/ollama serve"
tail -f /shared/$USER/ollama-serve.log

# --- model ----------------------------------------------------------
ollama list
ollama pull  qwen3.5:122b-a10b
ollama run   qwen3.5:122b-a10b --verbose "prompt"
ollama stop  qwen3.5:122b-a10b   # unload from VRAM, keep on disk
ollama rm    <model>             # delete weights from /shared

# capabilities decide what a model may be used for - check before wiring
curl -s http://localhost:11435/api/show -d '{"model":"qwen3.5:122b-a10b"}' \
  | python3 -c "import json,sys; print(json.load(sys.stdin)['capabilities'])"

# --- endpoint -------------------------------------------------------
curl http://localhost:11435/api/tags
curl http://localhost:11435/api/ps          # what is resident in VRAM
curl http://localhost:11435/api/generate \
  -d '{"model":"qwen3.5:122b-a10b","prompt":"...","stream":false}'
\end{lstlisting}

\subsection*{corpus, on the developer VM}

\begin{lstlisting}
# --- index ------------------------------------------------------------
./rag ingest                            # new uploads in ~/work/raw
./rag ingest /shared/eugen/brochures    # a folder elsewhere
./rag ingest --force                    # re-parse everything
./rag ingest --reset                    # rebuild from scratch
./rag status                            # what is indexed

# --- ask --------------------------------------------------------------
./rag ask    "..."                      # cited answer
./rag ask    "..." --fast               # skip the rerank call
./rag search "..."                      # passages only

# --- server -----------------------------------------------------------
tmux new -s corpus-mcp
./mcp-http --host 0.0.0.0               # MCP + JSON on :8765, bearer token
./mcp-http --show-client-config

# --- health check -----------------------------------------------------
# No token must return 401. A bare tools/list returns -32600 Missing
# session ID: initialize first, then reuse the mcp-session-id header.
# Full sequence in Section 9.5.
curl -s -o /dev/null -w '%{http_code}\n' -X POST \
  http://127.0.0.1:8765/mcp -H 'Content-Type: application/json' -d '{}'
\end{lstlisting}
